\documentclass[11pt]{article}

\usepackage[T1]{fontenc}
\usepackage{lmodern} % Use scalable vector fonts for microtype
\usepackage{amsmath,amssymb,amsthm,mathtools}
\usepackage{booktabs}
\usepackage{enumitem}
\usepackage[margin=1.15in]{geometry}
\usepackage{microtype}
\usepackage[hidelinks]{hyperref}

\newtheorem{theorem}{Theorem}[section]
\newtheorem{corollary}[theorem]{Corollary}
\newtheorem{proposition}[theorem]{Proposition}
\newtheorem{lemma}[theorem]{Lemma}
\newtheorem{remark}[theorem]{Remark}
\newtheorem{example}[theorem]{Example}

\newcommand{\C}{\mathbb{C}}
\newcommand{\Z}{\mathbb{Z}}
\newcommand{\gr}{\operatorname{gr}}
\newcommand{\Span}{\operatorname{span}}
\newcommand{\area}{\operatorname{area}}
\newcommand{\dinv}{\operatorname{dinv}}
\newcommand{\Dyck}{\operatorname{Dyck}}
\newcommand{\Prim}{\operatorname{Prim}}
\newcommand{\e}{\mathbf e}
\newcommand{\alg}{\mathrm{alg}}
\newcommand{\doi}[1]{\href{https://doi.org/#1}{doi:\,#1}}

\title{On the Middle Coefficients of Rational $q,t$-Catalan Polynomials}
\author{Graham Hawkes}
\date{August 2026}

\begin{document}

\maketitle

\begin{abstract}
Let $r,s$ be coprime positive integers and let
$g=(r-1)(s-1)/2$.  We prove that the part of the rational
$q,t$-Catalan polynomial $C_{r,s}(q,t)$ of total degree $g-\delta$
has a two-variable Schur expansion
\[
 C_{r,s;\delta}(q,t)
 =\sum_{a=0}^{\min\{\delta,\lfloor(g-\delta)/2\rfloor\}}
 c_{r,s;\delta,a}\,s_{(g-\delta-a,a)}(q,t),
 \qquad c_{r,s;\delta,a}\in\Z_{\geq 0}.
\]
Consequently, the coefficient of $q^jt^{g-\delta-j}$ is independent
of $j$ for
\[
 \delta\leq j\leq g-2\delta.
\]
The proof uses Ma's scalar rational-master character theorem for the
algebraic filtration on a finite-dimensional rational Cherednik
module.  That character is converted into the bigraded Hilbert series
of the maximal-ideal associated graded of the spherical module.  Each
fixed-deficit piece is an $\mathfrak{sl}_2$-module, and its primitive
multiplicities are controlled by a quotient generated in positive
deficit.  We also deduce a primitive Hodge-vanishing statement for the
Catalan variety of Galashin and Lam.  An appendix records the exact
version of Ma's theorem used here and gives a self-contained
normalization argument for the two global grading shifts.
\end{abstract}

\section{Introduction}

The cases $r=1$ or $s=1$ are immediate, since then the rational
Catalan polynomial is $1$.  We therefore assume that $r,s\geq2$ are
coprime and put
\[
 g=\frac{(r-1)(s-1)}2.
\]
In the area--dinv convention of the rational shuffle theorem
\cite{Mellit}, the rational $q,t$-Catalan polynomial is
\begin{equation}\label{eq:Catalan-definition}
 C_{r,s}(q,t)
 =\sum_{D\in\Dyck_{r,s}}q^{\area(D)}t^{\dinv(D)}.
\end{equation}
For $0\leq\delta\leq g$, define the fixed-deficit slice
\begin{equation}\label{eq:deficit-slice}
 C_{r,s;\delta}(q,t)
 =\sum_{A+B=g-\delta}
 [q^At^B]C_{r,s}(q,t)\,q^At^B.
\end{equation}
Thus every monomial in this slice has total degree
$N=g-\delta$.

Our main result gives more than a constant middle interval: it bounds
the support of the full two-variable Schur expansion.

\begin{theorem}[Schur-support bound]\label{thm:main}
For every $0\leq\delta\leq g$, there are nonnegative integers
$c_{r,s;\delta,a}$ such that, with $N=g-\delta$,
\begin{equation}\label{eq:main-Schur}
 C_{r,s;\delta}(q,t)
 =\sum_{a=0}^{\min\{\delta,\lfloor N/2\rfloor\}}
 c_{r,s;\delta,a}\,s_{(N-a,a)}(q,t).
\end{equation}
Equivalently, no constituent $s_{(N-a,a)}(q,t)$ occurs when
$a>\delta$.
\end{theorem}

Since
\begin{equation}\label{eq:two-variable-Schur}
 s_{(N-a,a)}(q,t)
 =\sum_{j=a}^{N-a}q^jt^{N-j},
\end{equation}
we obtain the following coefficient statement.

\begin{corollary}[Flat middle]\label{cor:flat-middle}
For every $0\leq\delta\leq g$, the coefficient
\[
 [q^jt^{g-\delta-j}]C_{r,s}(q,t)
\]
is independent of $j$ throughout
\begin{equation}\label{eq:flat-interval}
 \delta\leq j\leq g-2\delta.
\end{equation}
The interval is nonempty exactly when $3\delta\leq g$.
\end{corollary}

The analogous full statement in the classical case
$C_n(q,t)=C_{n+1,n}(q,t)$ was formulated by the author as
\cite[Conjecture~5.38]{HawkesDyck}.  Thus
Corollary~\ref{cor:flat-middle} proves that conjecture and gives its
rational extension.  Symmetric string decompositions and related
formulas were previously obtained in several special settings.  The
case with one parameter equal to $3$ has an explicit formula due to
Kaliszewski and Li \cite{KaliszewskiLi}; generalized $q,t$-Catalan
polynomials of length at most $4$ admit symmetric-chain decompositions
by Gorsky, Hawkes, Schilling, and Rainbolt \cite{GHSR}; and a
symmetric monomial-string formula for general rational parameters was
proposed in \cite{HawkesRational}.  The new point here is the uniform
primitive-support inequality $a\leq\delta$ for every coprime pair.

The proof is algebraic.  Let $L_{r/s}$ be the finite-dimensional
irreducible polynomial representation of the type-$A$ rational
Cherednik algebra, and let
\[
 R=\e L_{r/s}=L_{r/s}^{S_s}
\]
be its spherical part.  Ma's rational-master theorem identifies the
bigraded character of the algebraic filtration on $R$ with
$C_{r,s}$.  The Gorsky--Oblomkov--Rasmussen--Shende filtration then
converts this character into the Hilbert series of
\[
 G=\gr_{\mathfrak m}R,
\]
where $\mathfrak m=R_{>0}$.  The fixed-deficit subspace
\[
 G^{(\delta)}=\bigoplus_{\ell\geq0}G_{2\ell+\delta,\ell}
\]
is an actual finite-dimensional $\mathfrak{sl}_2$-module.  Its
raising operator is multiplication by $p_2/2$.  The multiplicity of a
string beginning in position $a$ is therefore a homogeneous dimension
of $G/(p_2)$.  Finally,
\[
 \C[z_3,\ldots,z_s]\twoheadrightarrow G/(p_2),
 \qquad \deg z_j=(j,1),
\]
and every generator $z_j$, $j\geq3$, contributes at least one unit of
deficit.  A monomial of maximal-ideal order $a$ must consequently
have deficit at least $a$.

The rational-master character framework was formulated by Gorsky and
Negu\c{t} \cite{GorskyNegut}.  The only input not presently available
in a refereed publication is the scalar rational-master character in
Ma's updated author manuscript \cite{MaCuspidal}.  Ma's filtration comparison is now
published \cite{MaFiltrations}.  Appendix~\ref{app:Ma} identifies the
exact manuscript version used here, records the logical dependencies,
and supplies a unit normalization that fixes both global grading
shifts.

The paper is organized as follows.  Section~\ref{sec:Cherednik-input}
states the Cherednik input and fixes filtration conventions.
Section~\ref{sec:m-adic} converts the character to the
$\mathfrak m$-adic associated graded.  Sections~\ref{sec:sl2} and
\ref{sec:support} construct the fixed-deficit $\mathfrak{sl}_2$-module
and prove the support bound.  Section~\ref{sec:consequences} gives
coefficient formulas and examples.  Section~\ref{sec:Hodge} records
the Hodge-theoretic corollary.

\section{The scalar Cherednik character}\label{sec:Cherednik-input}

Set $c=r/s$ and let $H_c(S_s,\mathfrak h)$ be the rational Cherednik
algebra of type $A_{s-1}$, where
\[
 \mathfrak h
 =\{(x_1,\ldots,x_s)\in\C^s:x_1+\cdots+x_s=0\}.
\]
Let $L_c$ be the unique finite-dimensional irreducible quotient of the
polynomial representation, let
\[
 \e=\frac1{s!}\sum_{w\in S_s}w,
\]
and put
\[
 R=\e L_c=L_c^{S_s}.
\]
Let $h$ be the Cartan element in the internal Cherednik
$\mathfrak{sl}_2$, and write
\[
 R=\bigoplus_{k\in\Z}R(k)
\]
for the $h$-weight decomposition.  We use the standard normalization
in which the constant polynomial has weight $-g$; a homogeneous
polynomial of ordinary degree $u$ then has weight
\begin{equation}\label{eq:Euler-degree}
 k=u-g.
\end{equation}

Let $F_i^{\alg}L_c$ be the algebraic filtration of
Gorsky--Oblomkov--Rasmussen--Shende \cite{GORS}, equivalently Ma's
algebraic/Kazhdan filtration, and set
\[
 F_i^{\alg}R=R\cap F_i^{\alg}L_c.
\]
Every filtration step is $S_s$-stable.  Since the symmetrizer is an
exact projection in characteristic zero,
\[
 \gr^{\alg}R\cong \e\,\gr^{\alg}L_c
\]
as a bigraded vector space.

The character theorem needed below is the following scalar
specialization of Ma's rational-master theorem.

\begin{theorem}[Ma's scalar rational-master character]
\label{thm:Ma-character}
For coprime positive integers $r,s$,
\begin{equation}\label{eq:Ma-character}
 \sum_{i,k}
 \dim\gr_i^{\alg}R(k)\,x^ky^i
 =C_{r,s}(xy,x^{-1}y),
\end{equation}
where
\[
 \gr_i^{\alg}R(k)
 =\frac{F_i^{\alg}R(k)}{F_{i-1}^{\alg}R(k)}.
\]
\end{theorem}

Appendix~\ref{app:Ma} explains the exact source and normalization of
Theorem~\ref{thm:Ma-character}.  The remainder of the argument uses
that theorem, published results of \cite{GORS}, and elementary
commutative algebra.

For reference, the three filtrations used in the proof are summarized
below.

\begin{center}
\begin{tabular}{@{}lll@{}}
\toprule
Notation & Direction & Meaning \\
\midrule
$\mathcal F^iL_c$ & decreasing & modified powers filtration of \cite{GORS} \\
$F_i^{\alg}L_c$ & increasing & $(\mathcal F^iL_c)^\perp$ for the Dunkl form \\
$\mathfrak m^\ell R$ & decreasing & powers of the homogeneous maximal ideal of $R$ \\
\bottomrule
\end{tabular}
\end{center}

\section{The maximal-ideal associated graded}\label{sec:m-adic}

The polynomial representation quotient has the form
\[
 L_c\cong\C[\mathfrak h]/I_c
\]
for a homogeneous $S_s$-stable ideal $I_c$.  Exactness of invariants
gives
\[
 R\cong
 \C[\mathfrak h]^{S_s}/(I_c\cap\C[\mathfrak h]^{S_s}).
\]
For
\[
 p_j=\sum_{i=1}^s x_i^j,
\]
we have $p_1=0$ on $\mathfrak h$, and
\[
 \C[\mathfrak h]^{S_s}=\C[p_2,p_3,\ldots,p_s].
\]
Thus
\begin{equation}\label{eq:R-presentation}
 R\cong\C[p_2,p_3,\ldots,p_s]/J
\end{equation}
for a homogeneous ideal $J$.

Let
\[
 \mathfrak m=R_{>0}
\]
and form the doubly graded algebra
\[
 G=\gr_{\mathfrak m}R
 =\bigoplus_{u,\ell}G_{u,\ell},
 \qquad
 G_{u,\ell}
 =\frac{(\mathfrak m^\ell)_u}{(\mathfrak m^{\ell+1})_u}.
\]
Here $u$ is ordinary polynomial degree and $\ell$ is
$\mathfrak m$-adic order.

Let $\mathfrak a\subseteq L_c$ be the ideal generated by the
positive-degree invariant polynomials.  Under $R=L_c^{S_s}$,
\[
 \mathfrak a=\mathfrak mL_c.
\]

\begin{lemma}\label{lem:intersection}
For every $j\geq0$,
\begin{equation}\label{eq:intersection}
 \mathfrak a^j\cap R=\mathfrak m^j.
\end{equation}
\end{lemma}

\begin{proof}
The inclusion $\mathfrak m^j\subseteq\mathfrak a^j\cap R$ is
immediate.  Conversely, write an element
$v\in\mathfrak m^jL_c\cap R$ as
\[
 v=\sum_\nu r_\nu w_\nu,
 \qquad r_\nu\in\mathfrak m^j,
 \quad w_\nu\in L_c.
\]
Let $\rho:L_c\to L_c^{S_s}=R$ be the Reynolds operator.  Since the
$r_\nu$ are invariant, $\rho$ is $R$-linear.  Hence
\[
 v=\rho(v)=\sum_\nu r_\nu\rho(w_\nu)\in\mathfrak m^j.
\]
\end{proof}

Gorsky--Oblomkov--Rasmussen--Shende define the modified decreasing
filtration
\begin{equation}\label{eq:GORS-decreasing}
 \mathcal F^iL_c
 =\sum_{j\geq0}
 \left(
 \mathfrak a^j\cap
 \bigoplus_{k<2j-i}L_c(k)
 \right).
\end{equation}
Their Proposition~4.5 shows that every $\mathcal F^iL_c$ is preserved
by the internal $\mathfrak{sl}_2$.

\begin{lemma}\label{lem:restricted-filtration}
At Euler weight $k$,
\begin{equation}\label{eq:restricted-filtration}
 \mathcal F^iR(k)
 =\mathfrak m^{\lfloor(i+k)/2\rfloor+1}\cap R(k),
\end{equation}
where $\mathfrak m^j=R$ for $j\leq0$.
\end{lemma}

\begin{proof}
Each summand in \eqref{eq:GORS-decreasing} is $S_s$-stable, so taking
invariants commutes with the sum.  At weight $k$, the summand indexed
by $j$ occurs precisely when
\[
 k<2j-i,
 \qquad\text{or equivalently}\qquad
 j>\frac{i+k}{2}.
\]
The powers of $\mathfrak m$ form a decreasing chain, so the sum is the
smallest occurring power,
$\mathfrak m^{\lfloor(i+k)/2\rfloor+1}$.  Apply
Lemma~\ref{lem:intersection}.
\end{proof}

The algebraic filtration is defined using the nondegenerate Dunkl
form:
\begin{equation}\label{eq:orthogonal-filtration}
 F_i^{\alg}L_c=(\mathcal F^iL_c)^\perp.
\end{equation}
The form is homogeneous and $S_s$-invariant.  Its restriction to $R$
is nondegenerate: for $v\in R$ and $w\in L_c$, invariance gives
$(v,w)=(v,\rho(w))$.  Thus a vector in $R$ orthogonal to $R$ is
orthogonal to all of $L_c$ and is zero.  The same averaging argument
gives
\[
 F_i^{\alg}R=(\mathcal F^iR)^{\perp_R}.
\]
Consequently,
\begin{equation}\label{eq:duality}
 \gr_i^{\alg}R(k)
 \cong
 \left(
 \frac{\mathcal F^{i-1}R(k)}{\mathcal F^iR(k)}
 \right)^*
\end{equation}
as vector spaces.  We use only equality of dimensions; no equivariant
identification with the dual is needed.

\begin{proposition}\label{prop:gr-comparison}
If $i+k$ is odd, then $\gr_i^{\alg}R(k)=0$.  If
$i+k=2\ell$, then
\begin{equation}\label{eq:gr-comparison}
 \dim\gr_i^{\alg}R(k)=\dim G_{g+k,\ell}.
\end{equation}
\end{proposition}

\begin{proof}
If $i+k=2\ell$, Lemma~\ref{lem:restricted-filtration} gives
\[
 \mathcal F^{i-1}R(k)=\mathfrak m^\ell\cap R(k),
 \qquad
 \mathcal F^iR(k)=\mathfrak m^{\ell+1}\cap R(k).
\]
By \eqref{eq:Euler-degree}, weight $k$ is ordinary degree $u=g+k$,
so the quotient in \eqref{eq:duality} is $G_{g+k,\ell}$.  If $i+k$
is odd, the powers appearing for $i-1$ and $i$ coincide.
\end{proof}

We can now rewrite the rational-master character entirely in terms of
$G$.

\begin{proposition}\label{prop:Catalan-G-character}
One has
\begin{equation}\label{eq:Catalan-G-character}
 C_{r,s}(q,t)
 =\sum_{u,\ell}\dim G_{u,\ell}\,
 q^\ell t^{g-u+\ell}.
\end{equation}
In particular, the deficit of the summand indexed by $(u,\ell)$ is
\begin{equation}\label{eq:deficit-u-ell}
 \delta=u-2\ell.
\end{equation}
Thus, if
\begin{equation}\label{eq:G-delta}
 G^{(\delta)}
 =\bigoplus_{\ell\geq0}G_{2\ell+\delta,\ell},
\end{equation}
then
\begin{equation}\label{eq:deficit-character}
 C_{r,s;\delta}(q,t)
 =\sum_{\ell\geq0}
 \dim G_{2\ell+\delta,\ell}\,
 q^\ell t^{g-\delta-\ell}.
\end{equation}
\end{proposition}

\begin{proof}
A monomial $q^At^B$ on the right side of
\eqref{eq:Ma-character} becomes
\[
 (xy)^A(x^{-1}y)^B=x^{A-B}y^{A+B}.
\]
Hence $k=A-B$ and $i=A+B$.  Proposition~\ref{prop:gr-comparison}
then gives
\[
 \ell=\frac{i+k}{2}=A,
 \qquad
 u=g+k=g+A-B.
\]
Solving for $B$ yields $B=g-u+\ell$, which proves
\eqref{eq:Catalan-G-character}.  The total $q,t$-degree is
$g-u+2\ell$, giving \eqref{eq:deficit-u-ell} and
\eqref{eq:deficit-character}.
\end{proof}

The absence of negative deficits is also immediate from
\eqref{eq:R-presentation}: every element of $\mathfrak m^\ell$ has
ordinary degree at least $2\ell$, because the invariant algebra has no
degree-one generator.

\section{Fixed deficit and the internal \texorpdfstring{$\mathfrak{sl}_2$}{sl2}}
\label{sec:sl2}

The internal Cherednik $\mathfrak{sl}_2$ has raising operator
\begin{equation}\label{eq:E-p2}
 E=\frac12\left(p_2-\frac1s p_1^2\right)=\frac12p_2,
\end{equation}
since $p_1=0$ on $\mathfrak h$, and $[h,E]=2E$.

Fix $\delta$ and write
\[
 N=g-\delta.
\]
For $G_{2\ell+\delta,\ell}$, the Euler weight is
\[
 k=(2\ell+\delta)-g=2\ell-N.
\]
The relation $i+k=2\ell$ then gives $i=N$.  Summing over all weights
produces an actual weight-graded identification
\begin{equation}\label{eq:fixed-deficit-identification}
 \frac{\mathcal F^{N-1}R}{\mathcal F^NR}
 \cong
 \bigoplus_{\ell\geq0}G_{2\ell+\delta,\ell}
 =G^{(\delta)}.
\end{equation}
The quotient on the left is $\mathfrak{sl}_2$-stable by
\cite[Proposition~4.5]{GORS}; we transport this action to
$G^{(\delta)}$.

The following compatibility is the key bridge between the
representation-theoretic and commutative-algebraic descriptions.

\begin{lemma}\label{lem:E-compatibility}
Under \eqref{eq:fixed-deficit-identification}, the raising operator
$E=p_2/2$ is multiplication by the initial form of $p_2/2$ in
$G=\gr_{\mathfrak m}R$.  In bidegrees,
\begin{equation}\label{eq:E-bidegree}
 E:G_{u,\ell}\longrightarrow G_{u+2,\ell+1}.
\end{equation}
\end{lemma}

\begin{proof}
At weight $k=2\ell-N$, the corresponding summand of the left side of
\eqref{eq:fixed-deficit-identification} is
\[
 \frac{\mathfrak m^\ell\cap R(k)}
      {\mathfrak m^{\ell+1}\cap R(k)}.
\]
Multiplication by $p_2$ raises ordinary degree by $2$, Euler weight by
$2$, and maximal-ideal order by at least $1$.  Modulo
$\mathfrak m^{\ell+2}$, the induced map is precisely multiplication by
the initial form of $p_2$.  Dividing by $2$ gives the action of $E$.
\end{proof}

Every finite-dimensional $\mathfrak{sl}_2$-module is a direct sum of
irreducibles.  An irreducible whose lowest Euler weight is
$-N+2a$ has weights
\[
 -N+2a,-N+2a+2,\ldots,N-2a.
\]
By \eqref{eq:deficit-character}, its contribution to the fixed-deficit
character is
\[
 q^at^{N-a}+q^{a+1}t^{N-a-1}+\cdots+q^{N-a}t^a
 =s_{(N-a,a)}(q,t).
\]
We have therefore proved the preliminary Schur expansion
\begin{equation}\label{eq:preliminary-Schur}
 C_{r,s;\delta}(q,t)
 =\sum_{a=0}^{\lfloor N/2\rfloor}
 c_{r,s;\delta,a}\,s_{(N-a,a)}(q,t),
 \qquad c_{r,s;\delta,a}\in\Z_{\geq0}.
\end{equation}
The remaining issue is the support bound $a\leq\delta$.

For a finite-dimensional $\mathfrak{sl}_2$-module $V$, the quotient
$V/EV$ contains exactly the lowest-weight line of each irreducible
summand.  Applied to $G^{(\delta)}$, this gives
\begin{equation}\label{eq:primitive-G-delta}
 c_{r,s;\delta,a}
 =\dim\left(G^{(\delta)}/EG^{(\delta)}\right)_{2a+\delta,a}.
\end{equation}
Since $E$ preserves deficit, this is the corresponding component of
the global quotient:
\begin{equation}\label{eq:primitive-global}
 c_{r,s;\delta,a}
 =\dim(G/EG)_{2a+\delta,a}.
\end{equation}

\section{The primitive-support bound}\label{sec:support}

The presentation \eqref{eq:R-presentation} is a filtered surjection
for the homogeneous maximal-ideal filtrations.  The image of
$(p_2,\ldots,p_s)^\ell$ is exactly $\mathfrak m^\ell$, so passage to
associated gradeds gives a bigraded surjection
\begin{equation}\label{eq:polynomial-to-G}
 \C[z_2,z_3,\ldots,z_s]\twoheadrightarrow G,
 \qquad \deg z_j=(j,1),
\end{equation}
where $z_j$ is the initial form of $p_j$.  By
Lemma~\ref{lem:E-compatibility}, the image of $z_2$ is $2E$.
Quotienting by it yields
\begin{equation}\label{eq:polynomial-to-primitive}
 \C[z_3,z_4,\ldots,z_s]\twoheadrightarrow G/EG.
\end{equation}

A monomial $z_3^{b_3}\cdots z_s^{b_s}$ has maximal-ideal order
\[
 a=\sum_{j=3}^s b_j
\]
and ordinary degree
\[
 u=\sum_{j=3}^s j b_j.
\]
Its deficit is therefore
\begin{equation}\label{eq:monomial-deficit}
 \delta=u-2a
 =\sum_{j=3}^s(j-2)b_j
 \geq\sum_{j=3}^s b_j=a.
\end{equation}
Thus the source of \eqref{eq:polynomial-to-primitive} has no component
of bidegree $(2a+\delta,a)$ when $a>\delta$.  A quotient cannot create
a missing homogeneous component, so
\[
 (G/EG)_{2a+\delta,a}=0
 \qquad(a>\delta).
\]
Together with \eqref{eq:primitive-global}, this proves
$c_{r,s;\delta,a}=0$ for $a>\delta$.  Substituting into
\eqref{eq:preliminary-Schur} proves
Theorem~\ref{thm:main}.

\begin{proof}[Proof of Corollary~\ref{cor:flat-middle}]
Let $N=g-\delta$.  Every occurring string in
\eqref{eq:main-Schur} begins at some $a\leq\delta$ and ends at
$N-a\geq N-\delta$.  Hence every such string contains
$q^jt^{N-j}$ exactly once whenever
\[
 \delta\leq j\leq N-\delta=g-2\delta.
\]
The coefficient there is $\sum_a c_{r,s;\delta,a}$, independent of
$j$.
\end{proof}

\section{Consequences and examples}\label{sec:consequences}

The Schur expansion gives the complete coefficient sequence in each
deficit.

\begin{corollary}\label{cor:coefficient-formula}
For $0\leq j\leq N=g-\delta$,
\begin{equation}\label{eq:coefficient-formula}
 [q^jt^{N-j}]C_{r,s;\delta}(q,t)
 =\sum_{0\leq a\leq\min\{\delta,j,N-j\}}
 c_{r,s;\delta,a}.
\end{equation}
Consequently, every fixed-deficit coefficient sequence is symmetric,
weakly increasing toward the center, constant on the interval
$[\delta,N-\delta]$ when that interval is nonempty, and weakly
decreasing afterward.  Moreover,
\begin{equation}\label{eq:coefficient-difference}
 [q^jt^{N-j}]C_{r,s;\delta}
 -[q^{j-1}t^{N-j+1}]C_{r,s;\delta}
 =c_{r,s;\delta,j}
\end{equation}
for $1\leq j\leq\min\{\delta,\lfloor N/2\rfloor\}$.
\end{corollary}

The free polynomial source in \eqref{eq:polynomial-to-primitive}
also gives a useful upper bound.

\begin{proposition}\label{prop:partition-bound}
For all $a,\delta\geq0$,
\begin{equation}\label{eq:partition-bound}
 c_{r,s;\delta,a}
 \leq
 [X^aY^\delta]
 \prod_{j=3}^s\frac1{1-XY^{j-2}}.
\end{equation}
In particular, nonvanishing implies
\begin{equation}\label{eq:two-sided-support}
 a\leq\delta\leq(s-2)a.
\end{equation}
If $s\geq3$ and $a=\delta$, the only possible source monomial is
$z_3^\delta$, so
\begin{equation}\label{eq:outer-multiplicity}
 c_{r,s;\delta,\delta}\leq1.
\end{equation}
For $s=2$, the source in \eqref{eq:polynomial-to-primitive} is
$\C$, and only $a=\delta=0$ occurs.
\end{proposition}

\begin{proof}
The coefficient on the right side of \eqref{eq:partition-bound}
counts monomials satisfying
\[
 \sum_{j=3}^s b_j=a,
 \qquad
 \sum_{j=3}^s(j-2)b_j=\delta.
\]
The bidegree component of a quotient cannot exceed the corresponding
component of its source.  The remaining assertions follow from
$1\leq j-2\leq s-2$.
\end{proof}

When the interval in Corollary~\ref{cor:flat-middle} is nonempty, its
common coefficient has the representation-theoretic interpretation
\begin{equation}\label{eq:common-middle}
 \sum_a c_{r,s;\delta,a}
 =\dim\bigl(G^{(\delta)}/EG^{(\delta)}\bigr).
\end{equation}

\begin{example}[The pair $(3,4)$]\label{ex:3-4}
Here $g=3$, and
\[
 C_{3,4}(q,t)
 =q^3+q^2t+qt^2+t^3+qt
 =s_{(3)}(q,t)+s_{(1,1)}(q,t).
\]
The primitive positions are $(\delta,a)=(0,0)$ and $(1,1)$.
\end{example}

\begin{example}[The pair $(3,5)$]\label{ex:3-5}
The explicit rank-word formula of Kaliszewski and Li
\cite{KaliszewskiLi} gives
\begin{align*}
 C_{3,5}(q,t)
 &=q^4+q^3t+q^2t^2+qt^3+t^4+q^2t+qt^2\\
 &=s_{(4)}(q,t)+s_{(2,1)}(q,t).
\end{align*}
Thus the deficit-zero slice is one string of length $5$, while the
deficit-one slice is one string of length $2$.  The latter begins at
$a=\delta=1$, illustrating the sharpness of the support bound.
\end{example}

\begin{remark}[A parameter equal to $3$]\label{rem:rank-three}
When the Cherednik denominator is $s=3$, the primitive source is
$\C[z_3]$.  Hence
\[
 c_{r,3;\delta,a}=0\quad\text{unless }a=\delta,
 \qquad
 c_{r,3;\delta,\delta}\leq1.
\]
Thus each nonzero deficit slice is a single Schur string.  This is the
algebraic shadow of the explicit one-string-per-skip formula in
\cite{KaliszewskiLi}.
\end{remark}

\section{Primitive Hodge vanishing}\label{sec:Hodge}

This section is a corollary of Theorem~\ref{thm:main}, not an input to
its proof.  Galashin and Lam define the Catalan variety
\[
 X_{r,r+s}^{\circ}=\Pi_{r,r+s}^{\circ}/T,
\]
a smooth affine complex variety of dimension
$(r-1)(s-1)=2g$.  Their Theorem~1.1 and equation~(1.11) identify its
mixed Hodge polynomial with the rational Catalan polynomial
\cite{GalashinLam}:
\begin{equation}\label{eq:Hodge-Catalan}
 C_{r,s}(q,t)
 =\sum_{K,P}
 q^{P-K/2}t^{g-K/2}
 \dim H^{K,(P,P)}(X_{r,r+s}^{\circ},\C).
\end{equation}

The variety satisfies curious Lefschetz.  There is a class
\[
 \gamma\in H^{2,(2,2)}(X_{r,r+s}^{\circ})
\]
whose cup product sends $(K,P)$ to $(K+2,P+2)$.  For a class in
$H^{K,(P,P)}$, set
\[
 \delta=K-P.
\]
Its total $q,t$-degree in \eqref{eq:Hodge-Catalan} is $g-\delta$, and
its $q$-exponent is
\[
 a=P-\frac K2=\frac{P-\delta}{2}.
\]
Equivalently,
\begin{equation}\label{eq:Hodge-change}
 P=2a+\delta,
 \qquad
 K=2a+2\delta.
\end{equation}
Cup product by $\gamma$ sends $a$ to $a+1$, exactly as the Cherednik
raising operator does.

Define the primitive-source quotient
\begin{equation}\label{eq:Hodge-primitive}
 \Prim_\gamma^{P,\delta}(X_{r,r+s}^{\circ})
 =
 \frac{H^{P+\delta,(P,P)}(X_{r,r+s}^{\circ})}
 {\gamma\,H^{P+\delta-2,(P-2,P-2)}(X_{r,r+s}^{\circ})},
\end{equation}
with a missing source group interpreted as zero.  By uniqueness of the
Lefschetz--Schur string decomposition,
\begin{equation}\label{eq:c-Hodge-primitive}
 c_{r,s;\delta,a}
 =\dim\Prim_\gamma^{2a+\delta,\delta}
 (X_{r,r+s}^{\circ}).
\end{equation}
No direct comparison map between this quotient and $G/EG$ is needed:
both sets of primitive dimensions are determined by the same Schur
expansion of $C_{r,s}(q,t)$.

Above the middle, curious Lefschetz makes cup product by $\gamma$
surjective from the preceding Hodge position, so the quotient
\eqref{eq:Hodge-primitive} vanishes for $P>g$.  In the source range
$P\leq g$, the support bound $a\leq\delta$ is, by
\eqref{eq:Hodge-change}, equivalent to $P\leq3\delta$.  We conclude:

\begin{corollary}[Primitive Hodge vanishing]\label{cor:Hodge-vanishing}
For every $P$ and $\delta$,
\begin{equation}\label{eq:Hodge-vanishing}
 \Prim_\gamma^{P,\delta}(X_{r,r+s}^{\circ})=0
 \qquad\text{whenever }P>3\delta.
\end{equation}
Equivalently, writing $K=P+\delta$, the primitive-source quotient in
$H^{K,(K-\delta,K-\delta)}$ vanishes whenever
\begin{equation}\label{eq:Hodge-vanishing-K}
 K>4\delta.
\end{equation}
\end{corollary}

\appendix

\section{The rational-master input and its normalization}
\label{app:Ma}

The proof uses the scalar character in
Theorem~\ref{thm:Ma-character}.  Because the relevant rational-master
manuscript is currently an author version, we record the precise
source and the normalization needed here.

\subsection{Version and logical dependencies}

The arXiv record is \href{https://arxiv.org/abs/2407.00971}{arXiv:2407.00971}.
The version checked for this paper is the $42$-page author PDF linked
from Ma's research webpage as ``more updated than arxiv version,''
with PDF generation date 6 October 2025:
\begin{center}
\url{https://xinchma.github.io/pdfs/cuspidal_catalan.pdf}
\end{center}
Its SHA-256 digest is
\begin{center}
\texttt{1a4bc0ca37d9c7121274248d2792ea6800fd25625af224d2b3e91664d78d9e2d}.
\end{center}
The distinction from the older arXiv version matters because the
updated manuscript uses the rational-Catalan normalization required in
\eqref{eq:Ma-character}.

In the checked version, the route to the scalar theorem is as follows.
The rational shuffle theorem is stated as Theorem~5.4, with
its equation~(76) identifying the scalar spherical coefficient with
$C_{r,s}(q,t)$.  Theorem~5.5(b) computes the spherical character for
the Hodge filtration.  Definition~5.1 converts the Hodge filtration
and Euler grading into the Kazhdan grading.  Propositions~5.9 and 5.10,
summarized in Corollary~1.3, transfer the result to the
inductive/algebraic filtration for all positive coprime slopes; the
all-slope passage uses the slope symmetry recorded in
\cite[Corollary~1.1]{GorskyArc}.  The filtration identification is also
supported by the published paper
\cite{MaFiltrations}, especially Lemma~3.9 and Theorem~3.16.

The stronger sheaf-level identification in Conjecture~4.6 of the
author manuscript is not used.  The character calculation proceeds
through the equality of equivariant $K$-classes in Proposition~4.9,
together with localization and the dimension comparison in
Theorem~5.5.

\subsection{Fixing both global grading shifts}

The updated manuscript's discussion of extremal Euler degrees fixes
the Euler normalization.  Since the character has two gradings, we
record an additional normalization by the spherical unit; this also
fixes the filtration shift.

Suppose that after the Hodge-to-Kazhdan conversion the scalar character
were determined only up to a global factor $x^\alpha y^\beta$.  The
spherical module $R$ has a unique lowest Euler-weight line, spanned by
the constant polynomial.  Its Euler weight is $-g$ by the standard
normalization \cite[Theorem~1.10]{BEG}; see also
\cite[Section~2.2.3]{GORS}.  Its algebraic filtration degree is $g$:
in the notation of Proposition~\ref{prop:gr-comparison}, the unit is
$G_{0,0}$, so $k=-g$, $\ell=0$, and $i+k=2\ell$ gives $i=g$.
Therefore the unit contributes
\[
 x^{-g}y^g.
\]

On the Catalan side, there is a unique area-zero path, and its monomial
is $t^g$.  Under the substitution $(q,t)=(xy,x^{-1}y)$, this monomial
is again
\[
 (x^{-1}y)^g=x^{-g}y^g.
\]
Thus the unique lowest-weight term agrees on both sides, forcing
\[
 \alpha=\beta=0.
\]
This completes the normalization of
Theorem~\ref{thm:Ma-character} in the form used in the present paper.

\begin{thebibliography}{99}

\bibitem{BEG}
Y.~Berest, P.~Etingof, and V.~Ginzburg,
\emph{Finite-dimensional representations of rational Cherednik algebras},
Int. Math. Res. Not. 2003, no.~19, 1053--1088,
\doi{10.1155/S1073792803210205}.

\bibitem{GalashinLam}
P.~Galashin and T.~Lam,
\emph{Positroids, knots, and $q,t$-Catalan numbers},
Duke Math. J. 173 (2024), no.~11, 2117--2195,
\doi{10.1215/00127094-2023-0049}.

\bibitem{GorskyArc}
E.~Gorsky,
\emph{Arc spaces and DAHA representations},
Selecta Math. (N.S.) 19 (2013), no.~1, 125--140,
\doi{10.1007/s00029-012-0096-0}.

\bibitem{GHSR}
E.~Gorsky, G.~Hawkes, A.~Schilling, and J.~Rainbolt,
\emph{Generalized $q,t$-Catalan numbers},
Algebraic Combinatorics 3 (2020), no.~4, 855--886,
\doi{10.5802/alco.120}.

\bibitem{GORS}
E.~Gorsky, A.~Oblomkov, J.~Rasmussen, and V.~Shende,
\emph{Torus knots and the rational DAHA},
Duke Math. J. 163 (2014), no.~14, 2709--2794,
\doi{10.1215/00127094-2827126}.

\bibitem{GorskyNegut}
E.~Gorsky and A.~Negu\c{t},
\emph{Refined knot invariants and Hilbert schemes},
J. Math. Pures Appl. 104 (2015), no.~3, 403--435,
\doi{10.1016/j.matpur.2015.03.003}.

\bibitem{HawkesRational}
G.~Hawkes,
\emph{A conjectured formula for the rational $q,t$-Catalan polynomial},
Ann. Comb. 28 (2024), 749--795,
\doi{10.1007/s00026-023-00662-2}.

\bibitem{HawkesDyck}
G.~Hawkes,
\emph{Dyck symmetric functions and applications to $q,t$-Catalan polynomials},
arXiv:2605.13003v1 (2026).

\bibitem{KaliszewskiLi}
R.~Kaliszewski and H.~Li,
\emph{The $(m,n)$-rational $q,t$-Catalan polynomials for $m=3$ and their
$q,t$-symmetry},
Discrete Math. Theor. Comput. Sci. Proc., FPSAC 2015, 769--780,
\doi{10.46298/dmtcs.2500}.

\bibitem{MaCuspidal}
X.~Ma,
\emph{From Cherednik algebras to knot homology via cuspidal
$\mathcal D$-modules},
arXiv:2407.00971 (2024); updated author manuscript, 42 pages,
PDF generated 6 October 2025.

\bibitem{MaFiltrations}
X.~Ma,
\emph{Rational Cherednik algebras and knot invariants: filtrations},
Int. Math. Res. Not. 2026, no.~10, rnag089,
\doi{10.1093/imrn/rnag089}.

\bibitem{Mellit}
A.~Mellit,
\emph{Toric braids and $(m,n)$-parking functions},
Duke Math. J. 170 (2021), no.~18, 4123--4169,
\doi{10.1215/00127094-2021-0011}.

\end{thebibliography}

\end{document}
